\section*{Identificación y gestión de fuentes radiactivas}

Cabe señalar que las fuentes radiactivas se clasifican en fuentes cerradas o selladas y fuentes abiertas o no selladas; las primeras son aquellas encapsuladas bajo una envoltura con suficiente resistencia mecánica para impedir que se establezca contacto con el radionúclido o que la sustancia radiactiva se disperse en las condiciones previsibles de utilización y desgaste. Una fuente abierta es toda fuente que no está sellada y que, en las condiciones normales de uso, puede producir contaminación, es decir, existe el riesgo de que parte del material radiactivo tenga contacto físico no deseado con cuerpos o materiales del medio que la rodea.

Respecto a los riesgos antes mencionados, puede decirse que una fuente cerrada, durante su uso normal, solo produce riesgo de irradiación externa mientras no pierda hermeticidad evolutiva, mientras que una fuente abierta podría dar lugar a irradiación externa e interna, pues la posibilidad de una contaminación por ingestión o inhalación es mucho mayor.

\chapter{Medición de radiación}

Dentro de las finalidades de la medición de la radiación está la vigilancia individual, la cual se entiende como las actividades cuyo objetivo es estimar el equivalente de dosis y la incorporación de radionúclidos en las personas ocupacionalmente expuestas.

La vigilancia radiológica de los trabajadores ocupacionalmente expuestos a las radiaciones requiere de un sistema completo de medición, evaluación y registro de todas las exposiciones a las radiaciones que pueda sufrir cada individuo.

Una de las características que hace que la radiación sea peligrosa es que no es posible detectarla con nuestros cinco sentidos, por lo que requerimos de algún dispositivo especial o instrumento que lo haga por nosotros.

Para esto, se cuenta con varios tipos de detectores a los que podríamos clasificar en:

\begin{itemize}
    \item \textbf{Dispositivos de vigilancia personal:} Tienen como propósito el llevar un registro de la dosis recibida por el trabajador o bien alertar en caso de que el nivel de radiación sea muy alto. Entre estos tenemos a los dosímetros de bolsillo, dosímetros de película o placas monitoras.
    \item \textbf{Instrumentos de monitoreo:} Se emplean como instrumentos de reconocimiento a fin de evitar que los niveles de radiación sean mayores que los permitidos por la ley.
\end{itemize}

Para evaluar los equivalentes de dosis recibidos se emplean dosímetros individuales, los cuales deben ser llevados siempre.
